% Template for analytic proof documents in analytical/<result>/.
%
% Copy this file to analytical/<result>/DRAFT.tex (and later PROOF.tex), then
% replace the example content. The rules below are the project standard; see
% analytical/README.md. Build with `latexmk -pdf DRAFT.tex` in the result
% directory and keep the log free of errors, undefined references and
% overfull boxes. Do not commit build output (.aux, .log, .pdf, ...).
%
% Rules:
% - Default article layout. Use only the packages loaded below; do not add
%   geometry or similar layout packages. Large tables may additionally use
%   array or longtable.
% - Inline math is \( ... \), never $ ... $. Write math as LaTeX, not as
%   Unicode symbols or code spans.
% - Every display is numbered: equation, align, gather, or equation with an
%   inner aligned/split. No \[ ... \], $$, starred environments, \nonumber,
%   \notag or \tag.
% - Statements use the theorem environments below and proofs use the proof
%   environment; never typeset "Lemma 3." or a QED box by hand.
% - Every internal cross-reference is \ref or \eqref to a \label. Never write
%   a number ("Lemma 3", "Section 2", "equation (4)") or a positional phrase
%   ("the display above", "the following table") in place of a reference.
%   Citations of other papers' numbering ("[x, Section II]") are fine.
% - Label prefixes: sec:, thm:, lem:, prop:, cor:, def:, rem:, eq:, tab:.
% - Lean declarations, modules and repository paths go in \texttt, with
%   underscores escaped.
% - Keep source lines at most 100 characters.

\documentclass{article}

\usepackage[T1]{fontenc}
\usepackage{lmodern}
\usepackage{amsmath}
\usepackage{amssymb}
\usepackage{amsthm}
\usepackage{mathtools}
\usepackage[hidelinks]{hyperref}

\theoremstyle{plain}
\newtheorem{theorem}{Theorem}
\newtheorem{lemma}{Lemma}
\newtheorem{proposition}{Proposition}
\newtheorem{corollary}{Corollary}
\theoremstyle{definition}
\newtheorem{definition}{Definition}
\theoremstyle{remark}
\newtheorem{remark}{Remark}

% Optional: number equations within sections, e.g. (2.3), for long documents.
% \numberwithin{equation}{section}

% Shared notation; delete what the document does not use.
\DeclareMathOperator{\wt}{wt}
\DeclareMathOperator{\supp}{supp}
\DeclareMathOperator{\rank}{rank}
\DeclareMathOperator{\Span}{span}
\DeclarePairedDelimiter{\abs}{\lvert}{\rvert}
\newcommand{\Ftwo}{\mathbb{F}_2}

\title{Title stating the proved result}
\author{}
\date{}

\begin{document}
\maketitle

% Grade line. Drafts stay PROPOSITION until review and Lean obligations close;
% a final PROOF.tex contains only THEOREM-grade claims (PROGRESS.md, "Claim
% grades"). State what has and has not been checked.
\textbf{Grade: PROPOSITION.} Complete hand proof; independent review and
guarded Lean replay are pending.

\begin{theorem}[main result]\label{thm:main}
  Let \(m\) be a positive integer. State every hypothesis and the complete
  conclusion here, for example
  \begin{equation}\label{eq:main-bound}
    \wt(x)\ge m \quad\text{for every nontrivial logical } x.
  \end{equation}
\end{theorem}

Summarize the strategy and point to its parts by label: the definitions are
in Section~\ref{sec:definitions}, the key estimate is
Lemma~\ref{lem:example}, and Section~\ref{sec:proof} proves
Theorem~\ref{thm:main}.

\section{Definitions}\label{sec:definitions}

\begin{definition}\label{def:weight}
  For a binary vector \(x\in\Ftwo^n\), the weight \(\wt(x)\) is the number
  of nonzero coordinates, and \(\supp(x)\) is the set of those coordinates.
\end{definition}

\section{Lemmas}\label{sec:lemmas}

\begin{lemma}[short descriptive title]\label{lem:example}
  For all \(x,y\in\Ftwo^n\),
  \begin{equation}\label{eq:triangle}
    \wt(x+y)\le\wt(x)+\wt(y).
  \end{equation}
\end{lemma}

\begin{proof}
  A coordinate of \(x+y\) is nonzero only if the same coordinate of \(x\) or
  of \(y\) is nonzero, so
  \begin{align}
    \supp(x+y)&\subseteq\supp(x)\cup\supp(y),\label{eq:support}\\
    \abs{\supp(x)\cup\supp(y)}&\le\abs{\supp(x)}+\abs{\supp(y)}.
  \end{align}
  Combining these with Definition~\ref{def:weight} gives~\eqref{eq:triangle}.
\end{proof}

Tables that the text cites are floats with a caption and a label, as in
Table~\ref{tab:example}.

\begin{table}[htbp]
  \centering
  \begin{tabular}{rrr}
    \(m\) & \(n\) & bound \\
    \hline
    \(3\) & \(10\) & \(3\) \\
    \(5\) & \(26\) & \(5\) \\
  \end{tabular}
  \caption{Example parameters.}\label{tab:example}
\end{table}

\section{Proof of the main result}\label{sec:proof}

Prove Theorem~\ref{thm:main} here, citing earlier results by label, for
example the inclusion~\eqref{eq:support} and Lemma~\ref{lem:example}. This
establishes~\eqref{eq:main-bound}.

\section{Lean correspondence}\label{sec:lean}

% Long \texttt names cannot hyphenate; sloppypar avoids overfull lines.
\begin{sloppypar}
Name the Lean modules and declarations that formalize each statement, for
example \texttt{QldpcVerification.Example.main\_bound} in
\texttt{lean/QldpcVerification/Example.lean}. Say which statements are not
formalized.
\end{sloppypar}

\section{Attribution and scope}\label{sec:attribution}

Cite prior work with \verb|\href| for a single link, for example
\href{https://arxiv.org/abs/2504.08887}{arXiv:2504.08887}, or with a
\texttt{references.bib} file in the result directory, \verb|\cite|, and the
two commented lines below. State what is new and make no unsupported
priority claim.

% \bibliographystyle{amsplain}
% \bibliography{references}

\end{document}
